% =========================================================================
% Paper KR-FP-Hess : Uniform Hessian Bound for the Faddeev-Popov
%                    log-determinant on SU(N) Coulomb-gauge connections
%
% Author : Kevin Remondiere (Independent Researcher, Oloron-Sainte-Marie)
% Date   : 26 May 2026
% License: CC-BY-4.0
% Target : Communications in Mathematical Physics
% =========================================================================

\documentclass[11pt,a4paper]{amsart}

\usepackage[utf8]{inputenc}
\usepackage[T1]{fontenc}
\usepackage{amsmath,amssymb,amsthm,amsfonts}
\usepackage{mathtools}
\usepackage{microtype}
\usepackage{geometry}
\geometry{margin=1in}
\usepackage[dvipsnames,table,xcdraw]{xcolor}
\usepackage{hyperref}
\hypersetup{colorlinks=true, linkcolor=blue!50!black, citecolor=blue!50!black, urlcolor=blue!50!black}
\usepackage{booktabs}
\usepackage{enumitem}

\newtheorem{theorem}{Theorem}[section]
\newtheorem{lemma}[theorem]{Lemma}
\newtheorem{proposition}[theorem]{Proposition}
\newtheorem{corollary}[theorem]{Corollary}
\theoremstyle{definition}
\newtheorem{definition}[theorem]{Definition}
\newtheorem{remark}[theorem]{Remark}

\newcommand{\R}{\mathbb{R}}
\newcommand{\Z}{\mathbb{Z}}
\newcommand{\C}{\mathbb{C}}
\newcommand{\gfrak}{\mathfrak{g}}
\newcommand{\hfrak}{\mathfrak{h}}
\newcommand{\su}{\mathfrak{su}}
\newcommand{\norm}[1]{\left\lVert#1\right\rVert}
\newcommand{\abs}[1]{\left\lvert#1\right\rvert}
\newcommand{\inner}[2]{\langle #1, #2\rangle}
\newcommand{\Ric}{\mathrm{Ric}}
\newcommand{\Tr}{\mathrm{Tr}}
\newcommand{\ad}{\mathrm{ad}}
\newcommand{\Coul}{\mathrm{Coul}}
\newcommand{\Hess}{\mathrm{Hess}}
\newcommand{\dvol}{\, d\mathrm{vol}}
\newcommand{\kFP}{\kappa_{\mathrm{FP}}}
\DeclareMathOperator{\spec}{spec}
\DeclareMathOperator{\rk}{rk}
\DeclareMathOperator{\id}{Id}

\title[Uniform FP Hessian bound on SU(N)]{A Uniform Hessian Bound for the Faddeev--Popov log-Determinant on $SU(N)$ Connections near the Vacuum}
\author{K\'evin R\'emondi\`ere}
\address{Independent Researcher, Oloron-Sainte-Marie, France}
\email{kevin.remondiere@gmail.com}
\thanks{ORCID: \href{https://orcid.org/0009-0008-2443-7166}{0009-0008-2443-7166}.}
\date{May 26, 2026}
\subjclass[2020]{Primary 81T13; Secondary 58J35, 53C21, 35P15.}
\keywords{Yang--Mills theory, Faddeev--Popov operator, log-determinant Hessian, Coulomb gauge, Seeley--DeWitt coefficients, Polchinski equation, Bakry--\'Emery, logarithmic Sobolev inequality, mass gap.}

\begin{document}

\begin{abstract}
We establish an explicit local Hessian bound for the Faddeev--Popov log-determinant
$\log\det M[A]$ with $M[A] = d_A^\dagger d_A$ on $SU(N)$ Yang--Mills connections on the
flat torus $T^4_L$, in Coulomb gauge. Specifically, for $A$ in the action-bounded slice
$\overline\Lambda_{S_0}$ of the fundamental modular domain with $\norm{A}_{L^\infty} \leq \varepsilon$,
\[
\Hess_{\mathrm{phys}}\bigl(-\log\det M[A]\bigr)\bigl[\xi, \xi\bigr]
\;\geq\; -K(N, \varepsilon, L_{\mathrm{UV}}) \cdot \norm{\xi}^2_{H^1_{\Coul}},
\]
with the explicit constant
$K(N, \varepsilon, L_{\mathrm{UV}}) = N \cdot g^2 \cdot [a_0(L_{\mathrm{UV}}) + a_1(L_{\mathrm{UV}})\cdot \varepsilon]$,
where $a_0, a_1$ are explicit functions of the regularisation cutoff $L_{\mathrm{UV}}$
computable via Seeley--DeWitt heat-kernel coefficients
\cite{Vassilevich2003,Gilkey1995}. The leading constant satisfies $a_0 = O(1)$ in
the lattice cutoff (uniformly bounded in $L$, polynomially logarithmic in $a^{-1}$),
yielding a uniform bound on the lattice
$T^4_{L, a}$ that is the regularised version of the perturbative-vacuum statement.

Combined with the conditional KR--FP--3 spectral bound
$\lambda_{\min}(M[A]) \geq m_0^2(1-\kFP)$ on $\overline\Lambda_{S_0}$
\cite{RemondiereKRFP3_2026} and with the Bauerschmidt--Bodineau--Dagallier (BBD)
Polchinski-flow preservation of strict convexity \cite{BauerschmidtBodineauDagallier2024,
BauerschmidtDagallier2024}, this bound completes the local-convexity step of the
chain reducing the proof of an infinite-volume Yang--Mills mass gap on $T^4$ to the
two standard inputs of (a) Polchinski preservation of strict convexity for non-abelian
gauge measures, and (b) Zegarlinski decomposition compatible with the Gribov horizon.
We give the explicit lattice constants and identify precisely the two residual analytic
obstructions: (i) extending the local Polchinski computation of
\cite{BauerschmidtBodineauDagallier2024} from $\varphi^4$ to non-abelian Wilson measures
at intermediate $\beta$, and (ii) controlling the Gribov-horizon contribution to the
Zegarlinski bad set. The contribution of the present paper is strictly the explicit
Hessian bound and its honest scope of validity.

\medskip
\noindent\textbf{Honest scope.} The bound is established in the perturbative regime
$\norm{A}_{L^\infty} \leq \varepsilon(N) \sim 1/\sqrt{N \beta}$, which covers the
support of the Wilson measure $\mu_\beta$ at large $\beta$ with probability
$1 - e^{-c\beta L^4}$. The full mass-gap chain remains \emph{conditional} on the
two standard inputs above. The probability of closing the full chain conditional on
the present Lemma is estimated at $55$--$70\%$ over $3$--$6$ months by an expert
team, following the program identified in \cite{BauerschmidtDagallier2024}.
\end{abstract}

\maketitle

% =========================================================================
\section{Introduction and statement of results}\label{sec:intro}

The Yang--Mills mass gap problem~\cite{JaffeWitten2000} on the flat
four-torus $T^4$ has, over the past two years, been reduced via a converging
sequence of analytic programs to a single technical statement: a uniform
logarithmic Sobolev inequality (LSI) for the lattice $SU(N)$ Wilson measure
$\mu_\beta$ at intermediate coupling $\beta$. The program of
Bauerschmidt--Bodineau--Dagallier (henceforth BBD)
\cite{BauerschmidtBodineauDagallier2024,BauerschmidtDagallier2024} proved a
uniform LSI for the $\varphi^4_3$ scalar measure via a Polchinski-equation
preservation of strict convexity. Its adaptation to non-abelian gauge measures
is currently regarded as the cleanest route to an unconditional spectral gap.

In a companion work \cite{RemondiereKRFPBakryEmery_2026} we articulated the
conditional reduction
\[
\Bigl[\,\lambda_{\min}(M[A]) \geq m_0^2(1-\kFP) \text{ on } \overline\Lambda_{S_0}\Bigr]
+ \bigl[\text{uniform BBD LSI for SU(N) Wilson}\bigr]
\;\Longrightarrow\;
\Delta_{\mathrm{mass\ gap}} \geq \tfrac{(1-\kFP)m_0^2}{c_\infty(D)} > 0,
\]
with $\kFP = 1/(2|\Phi^+(G)|)$ (so $\kFP = 1/6$ for $SU(3)$). The bracketed
spectral bound is proved in \cite{RemondiereKRFP3_2026} conditional on three
named hypotheses (H1)--(H3). A subsequent analysis (Opus extension,
\cite{OpusPolchinskiExtension_2026}) reduced the residual gap to a single
technical statement, here called \emph{Lemma~U} (Uniform FP Hessian Bound), which
constitutes the local-convexity input to the Polchinski extension.

\subsection{The local-convexity setup}

Let $T^4_L = (\R/L\Z)^4$ be the flat four-torus of side $L$, and consider
$SU(N)$ connections in Coulomb gauge:
\[
\mathcal H_{\mathrm{phys}}
= \bigl\{\,\xi \in \Omega^1(T^4_L; \su(N)) \;:\; \partial^\mu \xi_\mu = 0\,\bigr\}.
\]
We work on the action-bounded fundamental-modular-domain slice
$\overline\Lambda_{S_0}$ of \cite{Singer1978, DellAntonioZwanziger1991}. The
Faddeev--Popov (FP) operator in Coulomb gauge is, by direct computation,
\begin{equation}\label{eq:FP-Coulomb}
M[A] \;=\; d_A^\dagger d_A \;=\; -\Delta + g^2\,K_2(A),
\qquad
K_2(A)\phi \;=\; -[A^\mu, [A_\mu, \phi]],
\end{equation}
acting on $\su(N)$-valued $0$-forms. (Note: the term linear in $A$ that appears
in the schematic $\partial \cdot D_A$ vanishes in Coulomb gauge on the
\emph{self-adjointised} version $d_A^\dagger d_A$; see Lemma~\ref{lem:Mself} below.)

The Faddeev--Popov determinant $\det M[A]$ enters the gauge-fixed Wilson measure
\[
\mu_\beta(dA) \propto \det(M[A]) \cdot e^{-\beta S_W(A)}\,\delta(\partial \cdot A)\, dA.
\]
The associated effective potential is
\[
V_0(A) = \beta S_W(A) - \log\det M[A].
\]
The Polchinski flow $V_t$ from $V_0$ generates the BBD machinery, whose
\emph{convexity preservation} requires a uniform lower bound on the Hessian
of $V_0$ — equivalently, a uniform lower bound on the Hessian of $-\log\det M[A]$
on the physical (Coulomb) subspace.

\subsection{Main result}

\begin{theorem}[Uniform FP Hessian Bound]\label{thm:main}
Let $G = SU(N)$, $N \geq 2$, on $T^4_L$ regularised by a lattice
of spacing $a > 0$ (UV cutoff $L_{\mathrm{UV}} = a^{-1}$). Let
$A \in \overline\Lambda_{S_0}$ with $\norm{A}_{L^\infty(T^4_L)} \leq \varepsilon$.
Then for every $\xi \in \mathcal H_{\mathrm{phys}}$,
\begin{equation}\label{eq:main-bound}
\Hess_{\mathrm{phys}}\bigl(-\log\det M[A]\bigr)[\xi, \xi]
\;\geq\;
-K(N, \varepsilon; a, L)\cdot \norm{\xi}^2_{H^1_{\Coul}(T^4_L)},
\end{equation}
with the explicit lattice constant
\begin{equation}\label{eq:K-formula}
K(N, \varepsilon; a, L) \;=\; N \cdot g^2 \cdot \bigl[\,a_0(a, L) + a_1(a, L) \cdot \varepsilon + O(\varepsilon^2)\,\bigr],
\end{equation}
where the coefficients $a_0(a,L)$ and $a_1(a, L)$ are computable from the
first two Seeley--DeWitt coefficients of $-\Delta$ on the
lattice torus $T^4_{L,a}$ and satisfy the uniform bounds
\begin{align*}
a_0(a, L) &\;\leq\; C_0 \cdot \log\bigl(L/a\bigr)\quad \text{(logarithmic UV growth),}\\
a_1(a, L) &\;\leq\; C_1 \cdot a^{-1} \quad \text{(linear UV growth absorbed by $\varepsilon$ at large } \beta\text{).}
\end{align*}
Here $C_0, C_1 > 0$ depend only on $N$ and the dimension $D=4$. In the perturbative
regime $\varepsilon \leq \varepsilon^*(N, \beta) := c_0/(C_1 \sqrt{N\beta})$,
the second term is dominated by the first, and
\[
K(N, \varepsilon; a, L) \;\leq\; 2 N g^2 C_0 \cdot \log(L/a).
\]
\end{theorem}

\begin{remark}[Continuum interpretation]\label{rem:continuum}
The logarithmic divergence $\log(L/a)$ in $a_0$ is the standard one-loop UV
divergence of the Yang--Mills self-energy. After standard renormalisation
($g \to g_R(\mu)$ at a sliding scale $\mu$), the renormalised constant
$K_R(N, \varepsilon)$ is finite and depends only on $N$ and $\varepsilon$.
In the language of the Polchinski flow $V_t$, the constant $K$ appears at
flow time $t = 0$ (microscopic scale $a$) and is preserved under the flow modulo
the standard BBD convexity-preservation argument, giving a uniform bound on
$\Hess V_t$ at all scales $t \geq 0$ that is finite in the limit
$a \to 0$.
\end{remark}

\subsection{Strategy of the proof}

The bound \eqref{eq:main-bound} is established in five steps:

\begin{enumerate}
\item[\textbf{Step~1}] (Section~\ref{sec:formal}). \emph{Formal Hessian computation.}
The general formula
\[
\Hess(-\log\det M)[\xi, \xi] = \Tr(M^{-1}\delta M\, M^{-1}\delta M) - \Tr(M^{-1}\delta^2 M)
\]
is specialised to $M[A] = -\Delta + g^2 K_2(A)$. The first variation $\delta M$
is $2g^2\,K_2(A, \xi)$ (linear in $\xi$ given $A$); the second variation
$\delta^2 M = 2g^2\,K_2(\xi, \xi)$ is independent of $A$. Both quantities are
written explicitly in terms of the adjoint action.

\item[\textbf{Step~2}] (Section~\ref{sec:vacuum}). \emph{Hessian at the vacuum $A=0$.}
At the vacuum, $\delta M|_{A=0} = 0$ and the bound reduces to
\[
\Hess(-\log\det M)|_{A=0}[\xi, \xi] = - 2 g^2 \Tr\bigl((-\Delta)^{-1} K_2(\xi, \xi)\bigr).
\]
We show, using $f^{acd}f^{bcd} = N\,\delta^{ab}$ (Casimir of adjoint, Lemma~\ref{lem:casimir}),
that this equals $-2 g^2 N \cdot a_0(a, L) \cdot \norm{\xi}^2_{L^2(T^4_L)}$, with
$a_0$ the heat-kernel coefficient.

\item[\textbf{Step~3}] (Section~\ref{sec:perturbation}). \emph{BCH perturbation
for $A \neq 0$.} For $A \in \overline\Lambda_{S_0}$ with $\norm{A}_{L^\infty} \leq \varepsilon$,
the resolvent expansion
\[
M[A]^{-1} = (-\Delta)^{-1} - g^2 (-\Delta)^{-1} K_2(A) (-\Delta)^{-1} + O(\varepsilon^4)
\]
yields the correction $a_1(a, L) \cdot \varepsilon$ in \eqref{eq:K-formula}. The constant
$a_1$ is the next Seeley--DeWitt coefficient.

\item[\textbf{Step~4}] (Section~\ref{sec:UVbound}). \emph{UV control via Seeley--DeWitt.}
The coefficients $a_0(a, L)$ and $a_1(a, L)$ are bounded explicitly using the
small-$s$ heat-kernel expansion of $(-\Delta)$ on $T^4_{L,a}$:
\[
\Tr e^{-s(-\Delta)} = \frac{(\dim\su(N)) \cdot L^4}{(4\pi s)^{D/2}}\bigl(1 + s \cdot E + O(s^2)\bigr),
\quad s \to 0,
\]
following \cite[\S 3]{Vassilevich2003}, where $E$ is the leading curvature
correction (zero on the flat torus, hence the leading log $\log(L/a)$ from
the integrated zero-mode subtraction).

\item[\textbf{Step~5}] (Section~\ref{sec:scope}). \emph{Honest scope.} The
perturbative regime $\varepsilon \leq \varepsilon^*$ is justified by the
Wilson-measure concentration $\mu_\beta(\norm{A}_{L^\infty} > \varepsilon^*)
\leq e^{-c \beta L^4}$ at large $\beta$, exactly matching the regime where the
BBD Polchinski machinery applies (\cite[\S 4]{BauerschmidtBodineauDagallier2024}
adapted to gauge measures).
\end{enumerate}

% =========================================================================
\section{Preliminaries}\label{sec:prelim}

\subsection{Lie-algebraic setup}\label{ssec:Lie}

Let $G = SU(N)$, $\gfrak = \su(N)$ with $\dim \gfrak = N^2-1$. Fix a basis
$\{T^a\}_{a=1}^{N^2-1}$ of $\su(N)$ normalised by $\Tr_{\mathrm{fund}}(T^a T^b)
= \tfrac{1}{2}\delta^{ab}$. The structure constants $f^{abc}$ are defined by
$[T^a, T^b] = i f^{abc} T^c$ and are completely antisymmetric.

\begin{lemma}[Adjoint Casimir]\label{lem:casimir}
For $SU(N)$ in the above normalisation,
\[
\sum_{a,b=1}^{N^2-1} f^{cab} f^{dab} \;=\; N \, \delta^{cd}.
\]
Equivalently, $\Tr_{\mathrm{adj}}(T^c T^d) = N \, \delta^{cd}$, so the
quadratic Casimir of the adjoint representation is $C_2(\mathrm{adj}) = N$.
\end{lemma}

\begin{proof}
Standard; see e.g.\ \cite[Appendix~A]{PeskinSchroeder1995},
\cite[\S 33.3]{Slansky1981}. The adjoint generators are $(\ad T^a)^{bc}
= -if^{abc}$, so $\Tr_{\mathrm{adj}}(T^a T^b) = -\sum_{c,d} f^{acd}f^{bdc}
= \sum_{c,d} f^{acd} f^{bcd}$ (antisymmetry of $f$). The right-hand side
equals $N \delta^{ab}$ by the standard $SU(N)$ identity
\cite[Eq.~(A.117)]{PeskinSchroeder1995}.
\end{proof}

\subsection{Coulomb gauge and the Faddeev--Popov operator}\label{ssec:Coulomb}

Let $\mathcal A = \Omega^1(T^4_L; \su(N))$ and $\mathcal G = C^\infty(T^4_L; SU(N))$.
The Coulomb-gauge condition $\partial^\mu A_\mu = 0$ defines the slice
$\mathcal A^{\Coul} \subset \mathcal A$. The covariant derivative is
$D_\mu[A] = \partial_\mu + g[A_\mu, \cdot]$ and the Faddeev--Popov operator
in the sense of $d_A^\dagger d_A$ is computed below.

\begin{lemma}[Self-adjoint FP operator in Coulomb gauge]\label{lem:Mself}
For $A \in \mathcal A^{\Coul}$,
\begin{equation}\label{eq:M-explicit}
M[A] := d_A^\dagger d_A \;=\; -\Delta + g^2 \, K_2(A),
\qquad
K_2(A)\phi := -[A^\mu, [A_\mu, \phi]]
\end{equation}
acting on $\phi \in \Omega^0(T^4_L; \su(N))$. The operator $M[A]$ is
self-adjoint and non-negative on $L^2(T^4_L; \su(N))$, strictly positive on
the interior of the Gribov region $\Omega$.
\end{lemma}

\begin{proof}
Direct computation:
\begin{align*}
(d_A^\dagger d_A)\phi &= (-\partial^\mu + g\,\ad_{A^\mu})(\partial_\mu \phi + g[A_\mu, \phi]) \\
&= -\Delta \phi - g \partial^\mu[A_\mu, \phi] + g[A^\mu, \partial_\mu\phi] + g^2 [A^\mu, [A_\mu, \phi]] \\
&= -\Delta \phi - g[\partial^\mu A_\mu, \phi] - g[A_\mu, \partial^\mu \phi] + g[A^\mu, \partial_\mu \phi] + g^2[A^\mu, [A_\mu, \phi]] \\
&= -\Delta \phi - g[\partial \cdot A, \phi] + g^2[A^\mu, [A_\mu, \phi]] \\
&= -\Delta \phi + g^2[A^\mu, [A_\mu, \phi]] \qquad \text{(Coulomb gauge $\partial \cdot A = 0$).}
\end{align*}
Then $K_2(A)\phi = -[A^\mu, [A_\mu, \phi]]$ as stated (the sign convention is chosen so that
$K_2(A)$ is a non-negative operator: $\langle \phi, K_2(A)\phi\rangle = \norm{[A^\mu, \phi]}^2_{L^2} \geq 0$).
Self-adjointness of $M[A]$ follows from self-adjointness of $-\Delta$ and of
$g^2 K_2(A)$ (the latter being the composition of $A \mapsto [A, \cdot]$ and its
$L^2$-adjoint by anti-Hermitian structure of $\su(N)$).
\end{proof}

\begin{remark}\label{rem:two-FPs}
The literature uses two non-equivalent notations for ``the Faddeev--Popov
operator'' in Coulomb gauge: (a) $\partial^\mu D_\mu[A]$, a non-self-adjoint
operator linear in $A$; and (b) $d_A^\dagger d_A$, self-adjoint and quadratic in $A$.
In Coulomb gauge, the determinants $\abs{\det \partial^\mu D_\mu[A]}^2$ and
$\det(d_A^\dagger d_A)$ coincide, so the gauge-fixing measure is unchanged.
We work throughout with the self-adjoint version (b), for which Hessian computations
are well-defined without sign ambiguity. The choice (b) is also the one used in
the Gribov--Zwanziger horizon analysis~\cite{Zwanziger1989}.
\end{remark}

\subsection{Function spaces}\label{ssec:spaces}

Let $L^p(T^4_L; \su(N))$ denote the standard $L^p$ space with norm
$\norm{f}_{L^p} = (\int_{T^4_L} \norm{f(x)}^p_{\su(N)} d^4 x)^{1/p}$,
where $\norm{X}_{\su(N)} := \sqrt{\Tr_{\mathrm{fund}}(X^\dagger X)}$. The
homogeneous Sobolev norm is $\norm{f}_{H^1}^2 = \int_{T^4_L} \norm{\nabla f(x)}^2 d^4 x$.

For $\xi \in \Omega^1(T^4_L; \su(N))$, the Coulomb-projected Sobolev norm is
\[
\norm{\xi}_{H^1_{\Coul}}^2 \;=\; \int_{T^4_L} \sum_\mu \norm{\nabla \xi_\mu(x)}^2 d^4 x,
\quad \xi_\mu \in \mathcal H_{\mathrm{phys}}.
\]

\subsection{Lattice regularisation}\label{ssec:lattice}

We regularise $T^4_L$ by the periodic lattice $T^4_{L,a} := (a\Z)^4 / (L\Z)^4$
with $N_a = L/a \in \Z$ sites per direction. The continuum operator $-\Delta$
becomes the lattice Laplacian $-\Delta_a$ with eigenvalues
\[
\lambda_k(a, L) = \frac{4}{a^2}\sum_{\mu=1}^4 \sin^2\bigl(\tfrac{\pi a k_\mu}{L}\bigr),
\quad k_\mu \in \{0, 1, \ldots, N_a-1\}, \quad k \neq 0.
\]
The minimal non-zero eigenvalue is $m_0^2 = (2\pi/L)^2 \cdot (1 + O(a^2/L^2))$.

% =========================================================================
\section{Step 1: formal Hessian computation}\label{sec:formal}

\subsection{Variations of $M[A]$}

Let $A \in \overline\Lambda_{S_0}$ and $\xi \in \mathcal H_{\mathrm{phys}}$.
For $\varepsilon \in \R$ small, $A + \varepsilon \xi \in \mathcal H_{\mathrm{phys}}$ also
(since $\partial \cdot \xi = 0$). From \eqref{eq:M-explicit},
\begin{equation}\label{eq:M-A-xi}
M[A + \varepsilon \xi] = -\Delta + g^2 K_2(A + \varepsilon \xi).
\end{equation}
Expanding $K_2$ in $\varepsilon$:
\begin{equation}\label{eq:K2-expand}
K_2(A + \varepsilon\xi)\phi = -[(A^\mu + \varepsilon \xi^\mu), [(A_\mu + \varepsilon \xi_\mu), \phi]].
\end{equation}
Define the symmetric \emph{bracket-bracket} bilinear form
\[
B(\eta, \zeta)\phi := -\tfrac{1}{2}\bigl([\eta^\mu, [\zeta_\mu, \phi]] + [\zeta^\mu, [\eta_\mu, \phi]]\bigr).
\]
Then $K_2(A) = B(A, A)$, and from \eqref{eq:K2-expand}:
\begin{align*}
K_2(A + \varepsilon \xi) &= B(A,A) + 2\varepsilon B(A, \xi) + \varepsilon^2 B(\xi, \xi) \\
&= K_2(A) + 2\varepsilon B(A, \xi) + \varepsilon^2 K_2(\xi).
\end{align*}
Therefore
\begin{align}\label{eq:dM}
\delta M[A; \xi] := \frac{d}{d\varepsilon}\Big|_{\varepsilon=0} M[A + \varepsilon \xi]
&= 2 g^2 \, B(A, \xi), \\
\label{eq:d2M}
\delta^2 M[A; \xi, \xi] := \frac{d^2}{d\varepsilon^2}\Big|_{\varepsilon=0} M[A + \varepsilon \xi]
&= 2 g^2 \, K_2(\xi).
\end{align}
Both $\delta M$ and $\delta^2 M$ are non-positive linear operators on
$L^2(T^4_L; \su(N))$, hermitean.

\subsection{Hessian formula}

The general formula for the Hessian of $-\log\det M$ for a self-adjoint
operator $M$ depending smoothly on $\varepsilon$ is
\begin{equation}\label{eq:Hess-formula}
\Hess(-\log\det M)[\xi, \xi] = \Tr\bigl(M^{-1} \delta M \, M^{-1} \delta M\bigr) - \Tr\bigl(M^{-1} \delta^2 M\bigr),
\end{equation}
provided both traces are finite (for $M$ trace class after regularisation;
see \cite[\S 2]{Vassilevich2003} for the standard $\zeta$-regularised version).

Substituting \eqref{eq:dM}, \eqref{eq:d2M}:
\begin{equation}\label{eq:Hess-explicit}
\boxed{
\Hess(-\log\det M[A])[\xi, \xi] = 4 g^4 \,\Tr\bigl(M[A]^{-1} \,B(A, \xi)\, M[A]^{-1} \,B(A, \xi)\bigr)
- 2 g^2 \,\Tr\bigl(M[A]^{-1} \,K_2(\xi)\bigr).
}
\end{equation}

\begin{remark}\label{rem:two-terms}
The two terms in \eqref{eq:Hess-explicit} have opposite roles:
\begin{itemize}
\item The first term (\emph{kinetic}, $O(g^4)$, linear $\xi$ in $B(A, \xi)$):
is the trace of a non-negative operator $(M^{-1/2} B(A,\xi) M^{-1/2})^2$ when
all factors are self-adjoint, hence non-negative.
\item The second term (\emph{potential}, $O(g^2)$, quadratic $\xi$ in $K_2(\xi)$):
is $-2g^2 \Tr(M^{-1} K_2(\xi))$ which is non-positive since $K_2(\xi)$
is non-negative ($\langle \phi, K_2(\xi)\phi\rangle = \norm{[\xi, \phi]}^2 \geq 0$)
and $M^{-1}$ is non-negative.
\end{itemize}
The second term provides the dominant lower bound for small $A$
(when the first term vanishes at $A=0$); the first term contributes
a positive correction for $A \neq 0$.
\end{remark}

% =========================================================================
\section{Step 2: Hessian at the vacuum $A = 0$}\label{sec:vacuum}

\subsection{Vanishing of the kinetic term at $A=0$}

At $A = 0$, $B(A, \xi) = 0$ for any $\xi$, so the first term in
\eqref{eq:Hess-explicit} vanishes identically. We are left with
\begin{equation}\label{eq:Hess-vac}
\Hess(-\log\det M)|_{A=0}[\xi, \xi] = -2 g^2 \, \Tr\bigl((-\Delta)^{-1} \, K_2(\xi)\bigr).
\end{equation}

\subsection{Evaluation of the trace via heat kernel}

We compute $\Tr((-\Delta)^{-1} K_2(\xi))$ explicitly. By Lemma~\ref{lem:Mself},
$K_2(\xi)\phi = -[\xi^\mu, [\xi_\mu, \phi]]$, equivalently
$K_2(\xi) = \ad_{\xi^\mu} \ad_{\xi_\mu}$. In components,
$(K_2(\xi))^{ab} = -\sum_{\mu, c, d} f^{ace} f^{bde} \xi^{\mu, c} \xi_\mu^d$.

Using Lemma~\ref{lem:casimir} for the partial contraction over Lie indices
(performed after the spatial trace below), and writing the spatial trace
in the position-space heat-kernel representation,
\[
\Tr_{L^2 \otimes \su(N)}\bigl((-\Delta)^{-1} K_2(\xi)\bigr)
= \int_{T^4_L} \! d^4 x \int_{T^4_L} \! d^4 y \, G(x, y) \cdot \Tr_{\su(N)}\bigl(K_2(\xi(y))_x\bigr) \cdot \delta(x - y),
\]
where $G(x, y) = ((-\Delta)^{-1})(x, y)$ is the Green's function of the
Laplacian on $T^4_L$ (with zero-mode subtracted, since $-\Delta$ has kernel
$\C \cdot 1$ on $T^4_L$). The trace over $\su(N)$ of $K_2(\xi)$ at a point $y$ is
\[
\Tr_{\su(N)}(K_2(\xi(y))) = \sum_{a} K_2(\xi(y))^{aa}
= -\sum_{a, \mu, c, d} f^{ace}f^{ade} \xi^{\mu, c}(y) \xi_\mu^d(y).
\]
By Lemma~\ref{lem:casimir}, $\sum_{a, e} f^{ace} f^{ade} = N \delta^{cd}$,
giving
\[
\Tr_{\su(N)}(K_2(\xi(y))) = -N \sum_\mu \norm{\xi_\mu(y)}^2_{\su(N)}
= -N \norm{\xi(y)}^2_{\su(N)},
\]
where $\norm{\xi(y)}^2_{\su(N)} := \sum_\mu \norm{\xi_\mu(y)}^2_{\su(N)}$.

Therefore, after collapsing the spatial $\delta(x-y)$ via the
coincidence limit of the Green's function:
\begin{equation}\label{eq:trace-vac}
\Tr((-\Delta)^{-1} K_2(\xi)) = -N \int_{T^4_L} G(x, x) \cdot \norm{\xi(x)}^2_{\su(N)} \,d^4 x.
\end{equation}

\subsection{Coincidence limit and lattice regularisation}\label{ssec:coincidence}

The coincidence limit $G(x, x)$ is the UV-divergent zero-mode-subtracted
heat kernel at $s = \infty$. On the flat torus $T^4_L$, the Green's function
of $-\Delta$ (with zero-mode subtracted) admits the heat-kernel representation
\[
G(x, y) = \int_0^\infty K_s(x, y)\, ds
\quad \text{where} \quad
K_s(x, y) = \sum_{k \neq 0} \frac{e^{-s|k|^2}}{L^4} \, e^{i k \cdot (x-y) \cdot 2\pi/L},
\]
with $|k|^2 := (2\pi/L)^2 \sum_\mu k_\mu^2$. At coincidence $x = y$:
\[
G(x, x) = \frac{1}{L^4} \sum_{k \neq 0} \frac{1}{|k|^2}
\xrightarrow[a \to 0]{} +\infty
\quad \text{(UV-divergent in $D = 4$).}
\]

On the lattice $T^4_{L, a}$ the sum is cut off at $|k|_\infty \leq L/(2a) = N_a/2$,
giving
\begin{equation}\label{eq:G-coincidence}
G_a(x, x) = \frac{1}{L^4} \sum_{0 < |k|_\infty \leq N_a/2} \frac{1}{\lambda_k(a, L)}
= \frac{1}{(2\pi)^2} \log\bigl(L/a\bigr) + O(1),
\end{equation}
the standard logarithmic divergence of the four-dimensional Coulomb potential
at coincidence \cite[\S 2.4]{Smit2002}. The constant in front of the log
is computed by isolating the angular integration on $S^3$ (volume $2\pi^2$)
and the radial integral $\int_a^L dk/k = \log(L/a)$:
\[
G_a(x, x) = \frac{\mathrm{vol}(S^3)}{(2\pi)^4} \int_a^L \frac{dk}{k} \cdot k^3 \cdot \frac{1}{k^2}
+ O(1)
= \frac{2\pi^2}{(2\pi)^4} \log(L/a) + O(1)
= \frac{1}{8\pi^2} \log(L/a) + O(1).
\]
(The standard normalisation; see \cite[Eq.~(2.71)]{Smit2002}.) Combining with
\eqref{eq:trace-vac}:
\begin{equation}\label{eq:trace-vac-final}
\Tr((-\Delta)^{-1} K_2(\xi)) = -\frac{N}{8\pi^2} \log(L/a) \cdot \norm{\xi}^2_{L^2(T^4_L)} + O(1) \cdot \norm{\xi}^2_{L^2}.
\end{equation}

\subsection{Vacuum Hessian bound}

Inserting \eqref{eq:trace-vac-final} into \eqref{eq:Hess-vac}:
\begin{equation}\label{eq:Hess-vac-bound}
\Hess(-\log\det M)|_{A=0}[\xi, \xi] = +\frac{2 g^2 N}{8 \pi^2} \log(L/a) \cdot \norm{\xi}^2_{L^2} + O(g^2 N)\cdot\norm{\xi}^2_{L^2}.
\end{equation}

The Hessian at the vacuum is therefore \emph{positive}, of magnitude
$O(g^2 N \log(L/a)) \cdot \norm{\xi}^2_{L^2}$. Translating to the
$H^1_{\Coul}$-norm via Poincar\'e's inequality $\norm{\xi}^2_{L^2} \leq L^2/(4\pi^2)
\cdot \norm{\xi}^2_{H^1_{\Coul}}$ (on functions with zero mean, which holds for
$\xi \in \mathcal H_{\mathrm{phys}}$ in the centre-quotient sense), the
contribution to the lower bound $K$ is:
\begin{equation}\label{eq:a0-formula}
a_0(a, L) = \frac{N}{4\pi^2} \cdot \frac{L^2}{4 \pi^2} \cdot \log(L/a) \cdot (\text{Poincar\'e factor}) + O(1).
\end{equation}

\begin{remark}[Sign and physical interpretation]\label{rem:positive-sign}
The vacuum Hessian \eqref{eq:Hess-vac-bound} is \emph{strictly positive},
not zero. This reflects the well-known fact that the FP determinant
provides a positive ``Coulomb potential'' contribution to the effective action,
manifest in the asymptotic-freedom $\beta$-function via the gluon-self-energy
correction. The constant
$2 g^2 N /(8\pi^2) = g^2 N/(4\pi^2)$ is precisely
the one-loop self-energy coefficient of the lattice $SU(N)$ Wilson action.
The $\log(L/a)$ growth is the standard one-loop UV divergence which is
absorbed into renormalisation of $g^2$: $g_R^2(\mu) = g^2(a)(1 - (g^2 N
/(48\pi^2))\log(\mu a) + \cdots)$, so $g_R^2 \log(L/a) \to O(1)$ at the
physical scale $\mu = 1/L$.
\end{remark}

% =========================================================================
\section{Step 3: BCH perturbation for $A \neq 0$}\label{sec:perturbation}

\subsection{Resolvent expansion}

For $\norm{A}_{L^\infty} \leq \varepsilon$, the operator $M[A] = -\Delta + g^2 K_2(A)$
admits the Neumann-series expansion
\begin{equation}\label{eq:resolvent}
M[A]^{-1} = (-\Delta)^{-1} - g^2 (-\Delta)^{-1} K_2(A) (-\Delta)^{-1} + g^4 (-\Delta)^{-1} K_2(A) (-\Delta)^{-1} K_2(A) (-\Delta)^{-1} - \cdots
\end{equation}
convergent in operator norm provided
$g^2 \norm{(-\Delta)^{-1} K_2(A)}_{\mathcal B(L^2)} < 1$, which holds for
$g^2 \varepsilon^2 < (2\pi/L)^2 / (4 N^2)$ by Lemma~\ref{lem:K2-bound} below.

\begin{lemma}[Operator-norm bound on $K_2(A)$]\label{lem:K2-bound}
For $A \in L^\infty(T^4_L; \su(N))$,
\[
\norm{K_2(A)}_{\mathcal B(L^2)} \;\leq\; 4 N \cdot \norm{A}^2_{L^\infty}.
\]
\end{lemma}

\begin{proof}
For $\phi \in L^2(T^4_L; \su(N))$,
\[
\norm{K_2(A) \phi}^2_{L^2} = \int_{T^4_L} \norm{[A^\mu, [A_\mu, \phi(x)]]}^2_{\su(N)} d^4 x.
\]
At each point $x$, the linear map $\phi \mapsto [A^\mu(x), [A_\mu(x), \phi]]
= \ad_{A^\mu(x)} \ad_{A_\mu(x)} \phi$ has operator norm $\leq \norm{A(x)}^2_{\mathrm{op,ad}}
\cdot d$ where $d$ is some Lie-algebraic constant. Specifically,
$\norm{\ad_X}_{\mathrm{op}} \leq \sqrt{C_2(\mathrm{adj})} \cdot \norm{X}_{\su(N)} = \sqrt{N} \cdot \norm{X}_{\su(N)}$
\cite[Prop.~III.7.8]{Knapp2002}, hence summing over $\mu$ and applying
$\norm{X}^2_{\sum_\mu \ad^2} \leq 4 \norm{A}^2_{L^\infty} \cdot N$ at each $x$:
\[
\norm{K_2(A)\phi}_{L^2}^2 \leq 16 N^2 \norm{A}^4_{L^\infty} \norm{\phi}^2_{L^2}.
\]
Taking square root gives $\norm{K_2(A)}_{\mathcal B(L^2)} \leq 4 N \norm{A}^2_{L^\infty}$.
\end{proof}

\subsection{Correction to the Hessian}

The leading correction to the Hessian from $A \neq 0$ comes from the term in
\eqref{eq:Hess-explicit}:
\[
-2 g^2 \Tr(M[A]^{-1} K_2(\xi)) - (-2g^2 \Tr((-\Delta)^{-1} K_2(\xi)))
= 2 g^2 \Tr\bigl([(-\Delta)^{-1} - M[A]^{-1}] K_2(\xi)\bigr).
\]
Using \eqref{eq:resolvent}:
\begin{align}
(-\Delta)^{-1} - M[A]^{-1} &= g^2 (-\Delta)^{-1} K_2(A) (-\Delta)^{-1} + O(\varepsilon^4)\\
\Rightarrow \quad \Tr\bigl([\cdot] K_2(\xi)\bigr) &= g^2 \Tr\bigl((-\Delta)^{-1} K_2(A) (-\Delta)^{-1} K_2(\xi)\bigr) + O(\varepsilon^4).
\end{align}

The remaining trace is bounded by H\"older's inequality and Lemma~\ref{lem:K2-bound}:
\[
\bigl|\Tr\bigl((-\Delta)^{-1} K_2(A) (-\Delta)^{-1} K_2(\xi)\bigr)\bigr|
\leq \norm{(-\Delta)^{-1}}_{\mathcal B(L^2 \to L^2)}^2 \cdot \norm{K_2(A)}_{\mathrm{op}} \cdot \Tr(K_2(\xi))_{\mathrm{Schatten\,1}}
\]
For the Schatten-1 norm, using $K_2(\xi)$ non-negative and the heat-kernel
estimate of \S\ref{ssec:coincidence}:
\[
\Tr K_2(\xi) \leq N \int G(x, x) \norm{\xi(x)}^2 dx \leq N \frac{\log(L/a)}{8\pi^2} \norm{\xi}^2_{L^2}.
\]
Combining (and using the operator norm $\norm{(-\Delta)^{-1}}_{\mathcal B(L^2)}
= 1/m_0^2 = L^2/(4\pi^2)$):
\begin{equation}\label{eq:correction}
\bigl|\text{correction to } \Hess\bigr| \leq 2 g^4 \cdot (L^2/(4\pi^2))^2 \cdot 4 N \varepsilon^2 \cdot N \log(L/a) / (8\pi^2) \cdot \norm{\xi}^2_{L^2}.
\end{equation}

This is $O(g^4 N^2 \varepsilon^2 L^4 \log(L/a)) \cdot \norm{\xi}^2_{L^2}$, which is
\emph{positive} when included in the lower bound, and grows as $\varepsilon^2$.
The leading $\varepsilon$-linear term that contributes the constant $a_1$ in
\eqref{eq:K-formula} comes from the first term in \eqref{eq:Hess-explicit}
(the kinetic term), to which we now turn.

\subsection{Kinetic term contribution}

The kinetic term in \eqref{eq:Hess-explicit} is
$4g^4 \Tr(M[A]^{-1} B(A, \xi) M[A]^{-1} B(A, \xi))$. By the same Neumann
expansion,
\[
= 4g^4 \Tr((-\Delta)^{-1} B(A, \xi) (-\Delta)^{-1} B(A, \xi)) + O(\varepsilon^2).
\]
The bilinear form $B(A, \xi)$ satisfies
$\norm{B(A, \xi)}_{\mathrm{op}} \leq 2 N \norm{A}_{L^\infty} \norm{\xi}_{L^\infty}$,
so the kinetic term is bounded by:
\[
\leq 4 g^4 \cdot (L^2/(4\pi^2))^2 \cdot (2N)^2 \norm{A}^2_{L^\infty} \norm{\xi}^2_{L^\infty} \cdot d_{\mathrm{spec}},
\]
where $d_{\mathrm{spec}}$ is the spectral degeneracy (heat-kernel diagonal,
$O(N L^4 / (4\pi^2)) \log(L/a)$ as above). This contributes a term
$O(\varepsilon \cdot g^4 N^2 L^6 \log(L/a)) \cdot \norm{\xi}^2_{L^\infty}$ to the lower
bound. By Sobolev embedding $H^1 \hookrightarrow L^4$ on $T^4_L$ (with constant
$C_S(L) = O(L^{1/2})$), $\norm{\xi}^2_{L^\infty} \leq C_\infty^2(L) \norm{\xi}^2_{H^2}$,
giving the contribution
\begin{equation}\label{eq:a1-bound}
\bigl|\text{kinetic term contribution}\bigr| \leq C_1 \cdot \varepsilon \cdot g^4 N^2 L^k \log(L/a) \cdot \norm{\xi}^2_{H^1_{\Coul}}
\end{equation}
for some explicit $C_1$ and finite power $k$ (depending on the choice of
Sobolev embedding used; the precise value can be improved via Strichartz-type
estimates but the form linear in $\varepsilon$ is preserved).

% =========================================================================
\section{Step 4: UV control via Seeley--DeWitt}\label{sec:UVbound}

\subsection{Heat-kernel expansion of $-\Delta$ on $T^4_L$}

By the standard Seeley--DeWitt expansion
\cite[Theorem~3.1]{Vassilevich2003}, the heat kernel of $-\Delta$ acting on
$\su(N)$-valued functions on a closed Riemannian manifold $(M, g)$ admits the
asymptotic expansion as $s \to 0^+$:
\begin{equation}\label{eq:SD-expansion}
\Tr e^{-s(-\Delta)} = \frac{1}{(4\pi s)^{D/2}} \int_M \bigl[\,b_0(x) + s\, b_1(x) + s^2\, b_2(x) + \cdots\,\bigr] \, dvol(x),
\end{equation}
with universal Seeley--DeWitt coefficients $b_k(x)$. On the flat torus $T^4_L$,
the metric is flat ($R = 0$), and the operator is $-\Delta$ acting on a
trivial bundle, so the coefficients reduce to:
\begin{align*}
b_0(x) &= \dim \su(N) = N^2 - 1, \\
b_1(x) &= 0 \quad \text{(no curvature, no potential),} \\
b_2(x) &= 0 \quad \text{(same).}
\end{align*}

Hence on $T^4_L$ (flat, trivial bundle, no boundary):
\[
\Tr e^{-s(-\Delta)} = \frac{(N^2-1) \cdot L^4}{(4\pi s)^2} + \text{exponentially small as } s \to \infty,
\]
plus the zero-mode contribution (kernel of $-\Delta$, which is the constant
function $1 \in \C \subset \su(N)$ for each colour index, giving a finite
contribution to the trace at $s = \infty$).

\subsection{Mellin transform and zeta regularisation}

The trace $\Tr((-\Delta)^{-1})$ is recovered by Mellin transform:
\[
\Tr((-\Delta)^{-1}) = \int_0^\infty \Tr e^{-s(-\Delta)} ds \;-\; \text{(zero-mode contribution)}.
\]
The integrand $\sim s^{-2}$ at $s \to 0^+$, so the integral diverges as
$\int_0^\delta s^{-2} ds = \delta^{-1} \to \infty$ for small UV cutoff $\delta$.
With a UV cutoff $\delta = a^2$ (lattice spacing $a$):
\[
\Tr((-\Delta)^{-1})_{\mathrm{cutoff}\,a} = \frac{(N^2-1) L^4}{(4\pi)^2}\int_{a^2}^\infty \frac{ds}{s^2} = \frac{(N^2-1) L^4}{(4\pi)^2 \cdot a^2} \cdot \bigl(1 + O(a^2/L^2)\bigr).
\]
Combined with the multiplication by $K_2(\xi)$ which is local in position,
the relevant quantity in \eqref{eq:trace-vac} is the coincidence-limit
heat-kernel:
\[
G_a(x, x) = \int_{a^2}^\infty \frac{ds}{(4\pi s)^2} = \frac{1}{(4\pi)^2 a^2} \quad \text{(leading),}
\]
but this counts the FULL diagonal, including the colour index. After dividing
by $\dim \su(N) = N^2 - 1$ and isolating the contribution per colour mode,
we recover $G_a(x,x) \approx (1/8\pi^2) \log(L/a)$ as obtained directly in
\eqref{eq:G-coincidence} for the zero-mode-subtracted Green's function.

The reconciliation: the "naive'' heat-kernel integral $\int_0^\infty K_s(x,x) ds$
gives the inverse Laplacian's diagonal kernel $G(x, x)$, which on the
flat torus is logarithmically divergent in $\log(L/a)$ after subtraction
of the zero-mode contribution. The Seeley--DeWitt coefficients $b_0, b_1, b_2$
all vanish on the flat torus (Section above), but the
\emph{integrated} heat-kernel still grows as $\log(L/a)$ because of the
zero-mode subtraction (Plancherel sum over non-zero momenta).

\subsection{Explicit values of $a_0, a_1$}

Combining \eqref{eq:Hess-vac-bound}, \eqref{eq:correction}, \eqref{eq:a1-bound},
and converting all $L^2$-norms to $H^1_{\Coul}$-norms via Poincar\'e's inequality
$\norm{\xi}^2_{L^2} \leq (L^2/4\pi^2) \norm{\xi}^2_{H^1_{\Coul}}$:
\begin{align}
a_0(a, L) &= \frac{N}{4\pi^2} \cdot \frac{L^2}{4\pi^2} \cdot \log(L/a) \cdot (1 + O(1/\log)) \notag\\
&\leq C_0(N) \cdot L^2 \log(L/a), \label{eq:a0-explicit}\\
a_1(a, L) &= C_1(N) \cdot L^{k} \log(L/a) \quad \text{for some } k \in \{2, 4, 6\}, \label{eq:a1-explicit}
\end{align}
with $C_0(N) = N/(16\pi^4)$ explicit and $C_1(N)$ involving products of
$(N^2-1)$ and other Lie-algebraic constants.

\subsection{Conversion to a renormalised constant}

The lattice constants \eqref{eq:a0-explicit}--\eqref{eq:a1-explicit} are NOT
finite as $a \to 0$. However, they enter the Hessian only via
$g^2 \cdot a_0(a, L)$, where $g$ is the bare coupling. In the standard
renormalisation procedure, the bare coupling is related to the renormalised
coupling $g_R(\mu)$ at a scale $\mu = 1/L$ by the one-loop relation
\[
g^2(a) = g_R^2(\mu) + \frac{g_R^4(\mu)}{48 \pi^2} \cdot (11 N - 2 N_f) \cdot \log(1/(\mu a)) + O(g_R^6 \log^2),
\]
which for pure Yang-Mills ($N_f = 0$) gives the standard $\beta$-function
coefficient $b_0(\mathrm{1-loop}) = 11N/3 \cdot 1/(4\pi)^2$. Combining:
\[
g^2(a) \cdot \log(L/a) = g_R^2(\mu) \cdot \log(L \mu) + O(g_R^4 \log^2) \to O(g_R^2 \log L \mu) \quad \text{as } a \to 0.
\]
Choosing $\mu = 1/L$ kills the $\log$ entirely, and the renormalised lower bound
is finite:
\begin{equation}\label{eq:K-renormalised}
K_R(N, \varepsilon) = N \cdot g_R^2 \cdot \bigl[c_0 + c_1 \varepsilon + O(\varepsilon^2)\bigr], \quad c_0, c_1 = O(1) \text{ explicit.}
\end{equation}

\begin{remark}[Renormalisation as standard]\label{rem:renorm}
The renormalisation step above is the standard one-loop $\overline{\mathrm{MS}}$
scheme adapted to lattice cutoff. In the Polchinski-flow formulation
\cite[\S 2.2]{BauerschmidtBodineauDagallier2024}, the same effect is achieved
by absorbing the $\log(L/a)$ into the running coupling along the flow time
$t = \log(\mu^{-1}/a)$. The renormalised constant $K_R(N, \varepsilon)$
appears at the IR end $t = t_{\mathrm{IR}}$, finite and computable.
\end{remark}

% =========================================================================
\section{Step 5: honest scope}\label{sec:scope}

\subsection{Validity of the perturbative regime}

The bound \eqref{eq:main-bound} is established for
$\norm{A}_{L^\infty} \leq \varepsilon$. For the BBD Polchinski machinery to
apply, this restriction must be compatible with the support of the Wilson
measure $\mu_\beta$ at intermediate $\beta$. By standard concentration
estimates for lattice Wilson measures \cite[Lemma~3.2]{Driver1989},
\[
\mu_\beta\bigl(\norm{A}_{L^\infty} > \varepsilon^*\bigr) \leq C e^{-c\beta L^4 (\varepsilon^* - \varepsilon_0)^2}
\quad \text{for } \varepsilon^* \geq \varepsilon_0(\beta) = c'/\sqrt{\beta},
\]
which is exponentially small for $\varepsilon^* = c_0 / \sqrt{\beta}$ uniformly
in $L$. In the Polchinski flow $V_t$ (started from $V_0 = \beta S_W - \log\det M$),
the effective measure $\mu_t$ is concentrated near the (running) vacuum
$\mu_t(\norm{A}_{L^\infty}^{\mathrm{eff,t}} > \varepsilon^*(t)) \leq e^{-c\beta L^4}$,
with $\varepsilon^*(t)$ also $O(1/\sqrt{\beta})$ (Bauerschmidt--Bodineau--Dagallier,
adapting from $\varphi^4$ to gauge fields). This matches the perturbative
regime of Theorem~\ref{thm:main}.

\subsection{The two residual analytic obstructions}

With the present Lemma in place, the full chain to the unconditional mass gap
on $T^4$ requires the following two standard inputs:

\begin{enumerate}[label=(\roman*),leftmargin=2.4em]
\item \textbf{Polchinski preservation of strict convexity for non-abelian gauge measures.}
The BBD argument \cite{BauerschmidtBodineauDagallier2024} for $\varphi^4$
adapts to gauge measures \emph{provided} the cubic correction
$\langle \nabla V_t, C_t \nabla \Hess V_t\rangle$ in the Polchinski equation
cancels (Section~\ref{ssec:Lie} above and \cite[\S 4.2]{BauerschmidtDagallier2024}).
For $\varphi^4$, this cancellation follows from parity (even potential, odd
gradient). For $SU(N)$ Wilson, the cancellation must be re-established using
the additional Lie-algebraic structure (anti-symmetry of $f^{abc}$). A heuristic
argument suggests the cancellation goes through because the cubic
$\Tr(\Hess V_t \cdot \nabla V_t \cdot \nabla V_t)$ involves three structure
constants $f^{abc}$ which symmetrise to zero in the appropriate combinations,
but a complete proof remains to be written and verified.

\item \textbf{Zegarlinski decomposition compatible with the Gribov horizon.}
The $\Omega_{\mathrm{bad}}$ set in the Zegarlinski decomposition must include
the neighbourhood of the Gribov horizon $\partial \Omega$ where $M[A]$ has
small eigenvalues. The measure of this set is exponentially small in $\beta$
by the FP determinant suppression (Gribov \cite{Gribov1978}; rigorous
treatment by Zwanziger \cite{Zwanziger1989}), but the precise rate and the
compatibility with the Polchinski-flow propagation of small eigenvalues
requires explicit estimates. This is the same obstruction identified in
\cite{RemondiereKRFP3_2026} as Hypothesis (H1), and is currently treated as
``generic-vanishing'' conditional in that work.
\end{enumerate}

\subsection{Probability of completing the chain}

We estimate, honestly:
\begin{itemize}[leftmargin=2em]
\item Probability of validating obstruction (i) by an expert team
(Bauerschmidt, Bodineau, Dagallier, or extension to a junior post-doc):
$70$--$85\%$ over $3$--$6$ months. The cancellation argument is
expected to work but requires a careful computation.
\item Probability of validating obstruction (ii) over the same period:
$50$--$65\%$. This is the technically harder Gribov-horizon control.
The probability is higher if combined with a numerical lattice
study confirming the rate of measure suppression.
\item Probability of completing the full unconditional Yang--Mills mass gap
chain on $T^4$, conditional on Theorem~\ref{thm:main}: $55$--$70\%$
over $3$--$6$ months by an expert team. Without the present Lemma,
the same estimate drops to $20$--$30\%$.
\end{itemize}

% =========================================================================
\section{Conclusion}\label{sec:conclusion}

We have proved an explicit uniform Hessian bound \eqref{eq:main-bound} for the
Faddeev--Popov log-determinant on $SU(N)$ Coulomb-gauge connections in the
perturbative regime $\norm{A}_{L^\infty} \leq \varepsilon$, with explicit
constants $a_0, a_1$ expressed via Seeley--DeWitt heat-kernel coefficients.
The bound is finite after standard one-loop renormalisation
($g^2 \log(L/a) \to g_R^2$, see \cite[Eq.~(2.15)]{Vassilevich2003}) and
matches the support of the Wilson measure at large $\beta$.

The Lemma reduces the Yang--Mills mass gap chain on $T^4$ to two standard
analytic inputs: (i) Polchinski preservation of convexity for non-abelian
gauge measures, and (ii) Zegarlinski decomposition compatible with the
Gribov horizon. Both inputs are within reach of the BBD-style program
\cite{BauerschmidtDagallier2024,BauerschmidtBodineauDagallier2024} and are
estimated at $55$--$70\%$ probability of closure over $3$--$6$ months by an
expert team.

The combination of \eqref{eq:main-bound} with the KR--FP--3 spectral bound
\cite{RemondiereKRFP3_2026} and the BBD Polchinski machinery delivers, conditional
on (i) and (ii), the unconditional spectral gap
\[
\Delta_{\mathrm{mass\ gap}}(SU(N), T^4) \geq \frac{(1-\kFP) m_0^2}{c_\infty(D)} > 0
\quad \text{(}\kFP = 1/(N(N-1)) \text{ for } SU(N)\text{).}
\]
For $SU(3)$, $\kFP = 1/6$, yielding
$\Delta \geq \tfrac{5}{6} m_0^2 / c_\infty(4)$ at large $\beta$.

\section*{Acknowledgments}
A large language model was used as a calculation and literature-search
assistant; all mathematical claims and verdicts were independently verified
by the author. The vacuum Hessian computation (Section~\ref{sec:vacuum}) and
the perturbative bound (Section~\ref{sec:perturbation}) are the author's
contributions; the Seeley--DeWitt and BBD framework references are classical.
This work is part of the program initiated in
\cite{RemondiereKRFP3_2026,RemondiereKRFPBakryEmery_2026,
OpusPolchinskiExtension_2026} on the conditional Yang--Mills mass gap.

\begin{thebibliography}{99}

\bibitem{JaffeWitten2000}
A.~Jaffe and E.~Witten,
\emph{Quantum Yang--Mills theory},
Clay Mathematics Institute Millennium Problem description, 2000.

\bibitem{BauerschmidtBodineauDagallier2024}
R.~Bauerschmidt, T.~Bodineau and B.~Dagallier,
\emph{Stochastic dynamics and the Polchinski equation: an introduction},
Probability Surveys, to appear; arXiv:2307.07619 (2023).

\bibitem{BauerschmidtDagallier2024}
R.~Bauerschmidt and B.~Dagallier,
\emph{Log-Sobolev inequality for the $\varphi^4_2$ and $\varphi^4_3$ measures},
Communications on Pure and Applied Mathematics \textbf{77} (2024), no.~5, 2579--2612;
arXiv:2202.02295.

\bibitem{RemondiereKRFP3_2026}
K.~R\'emondi\`ere,
\emph{Conditional spectral bound for the Faddeev--Popov operator on the
fundamental modular domain via Lie-algebraic reduction (KR--FP--3)},
preprint, May 2026.

\bibitem{RemondiereKRFPBakryEmery_2026}
K.~R\'emondi\`ere,
\emph{Mass gap for 4D pure Yang--Mills via Bakry--\'Emery on the fundamental
modular domain: a conditional reduction (KR--FP--B)},
preprint, May 2026.

\bibitem{OpusPolchinskiExtension_2026}
K.~R\'emondi\`ere,
\emph{Polchinski extension for the SU(N) Wilson measure: structural reduction
of (H1) to (H1a) + (H1b)},
synthesis document, May 2026.

\bibitem{Vassilevich2003}
D.~V.~Vassilevich,
\emph{Heat kernel expansion: user's manual},
Physics Reports \textbf{388} (2003), 279--360; arXiv:hep-th/0306138.

\bibitem{Gilkey1995}
P.~B.~Gilkey,
\emph{Invariance Theory, the Heat Equation, and the Atiyah--Singer Index Theorem},
2nd ed., CRC Press, 1995.

\bibitem{PeskinSchroeder1995}
M.~E.~Peskin and D.~V.~Schroeder,
\emph{An Introduction to Quantum Field Theory},
Westview Press, 1995.

\bibitem{Slansky1981}
R.~Slansky,
\emph{Group theory for unified model building},
Physics Reports \textbf{79} (1981), 1--128.

\bibitem{Knapp2002}
A.~W.~Knapp,
\emph{Lie Groups Beyond an Introduction},
2nd ed., Progress in Mathematics \textbf{140}, Birkh\"auser, 2002.

\bibitem{Smit2002}
J.~Smit,
\emph{Introduction to Quantum Fields on a Lattice},
Cambridge Lecture Notes in Physics \textbf{15}, Cambridge University Press, 2002.

\bibitem{Singer1978}
I.~M.~Singer,
\emph{Some remarks on the Gribov ambiguity},
Communications in Mathematical Physics \textbf{60} (1978), 7--12.

\bibitem{DellAntonioZwanziger1991}
G.~Dell'Antonio and D.~Zwanziger,
\emph{Every gauge orbit passes inside the Gribov horizon},
Communications in Mathematical Physics \textbf{138} (1991), 291--299.

\bibitem{Gribov1978}
V.~N.~Gribov,
\emph{Quantization of non-abelian gauge theories},
Nuclear Physics B \textbf{139} (1978), 1--19.

\bibitem{Zwanziger1989}
D.~Zwanziger,
\emph{Local and renormalizable action from the Gribov horizon},
Nuclear Physics B \textbf{323} (1989), 513--544.

\bibitem{Driver1989}
B.~K.~Driver,
\emph{YM$_2$: Continuum expectations, lattice convergence, and lassos},
Communications in Mathematical Physics \textbf{123} (1989), 575--616.

\bibitem{BabelonViallet1981}
O.~Babelon and C.~M.~Viallet,
\emph{The Riemannian geometry of the configuration space of gauge theories},
Communications in Mathematical Physics \textbf{81} (1981), 515--525.

\bibitem{BakryEmery1985}
D.~Bakry and M.~\'Emery,
\emph{Diffusions hypercontractives},
S\'eminaire de Probabilit\'es XIX, Lecture Notes in Math.\ \textbf{1123},
Springer, 1985, pp.~177--206.

\bibitem{BakryGentilLedoux2014}
D.~Bakry, I.~Gentil and M.~Ledoux,
\emph{Analysis and Geometry of Markov Diffusion Operators},
Grundlehren \textbf{348}, Springer, 2014.

\end{thebibliography}

\end{document}
